\section{Introduction}

Knowledge graphs (KGs) \cite{kg} represent facts in the real world as triples, such as \textit{(Paris, capital\_of, France)}, facilitating the organization and scaling of information \cite{kg-info}, and have gained paramount importance in knowledge-based systems, such as retrieval-augmented generation and recommendation systems~\cite{recommend}.
% 
However, most large-scale KGs \cite{kg}, built using rule-based methods and statistics-based methods to extract web data, often contain noisy or incorrect triples. For instance, the widely used knowledge graph NELL \cite{survey} contains around 600K incorrect triples, which account for 26\% of the set of triples in NELL. Most knowledge graph-driven tasks assume all triples are correct \cite{survey}, overlooking the impact of errors, which significantly degrades the performance of downstream tasks. This highlights the urgent need for effective KG error detection methods.

Existing methods can be broadly categorized into two types based on their evidence utilization. First, explicit evidence-based methods extract paths or subgraphs from the knowledge graph that support the truth of triples and use these explicit evidences to predict the credibility of triples~\cite{shi}. Second, embedding-based methods predict the credibility of triples by embedding entities and relations into vector space and calculating the embedded representation of paths or subgraphs~\cite{kgttm,embedding2}. Moreover, when combined with contrastive learning and pre-trained models, their performance has shown significant improvement~\cite{kg-bert,cca,sesicl}.

However, these methods have two main limitations: (1) They evaluate triples from a single, fixed perspective, either relying on static structural patterns (e.g., static graph embeddings) or a unidirectional semantic method (e.g., text embeddings). This means they assume predefined, unchanging connections between entities, without considering the dynamic context of the triples being evaluated. As a result, they struggle to adapt when the graph structure or textual representations change.
%
(2) They lack transparency in the evaluation process, typically providing a single confidence score without sufficient context about how decisions are made. These two limitations hinder error identification and impede model improvement. 

To address the limitations of existing methods, we introduce a multi-agent framework, as shown in Figure \ref{framework}.
%
For each triple in the KG, we assign two agents to the head and tail entities. The \textit{Forward Agent} collects subgraphs with the entity as the head, and the \textit{Backward Agent} collects subgraphs with it as the tail, integrating multiple perspectives for error detection.
% 
Then, we process the collected subgraphs using a Graph Convolutional Network (GCN) for structural features and an LLM for semantic features. By concatenating the GCN and LLM embeddings, we combine structural and semantic information, leveraging the LLM's text generation capabilities for error detection.

Moreover, to address the lack of transparency during evaluation, we introduce a structured discussion and voting mechanism. The multiple agents independently evaluate each triple and then engage in multiple rounds of discussion until they reach a final decision. 
% 
This method enhances detection accuracy and increases transparency in the decision-making process. 
At the same time, our framework can utilize domain-specific knowledge graphs to train specialized agents for industrial scenarios, improving error detection accuracy. It addresses the limitations of existing LLMs, which struggle with domain-specific knowledge, thereby delivering significant industrial value.
% 
The main contribution are summarized as follows:

\begin{enumerate}
    \item To our knowledge, we are the first to introduce a multi-agent framework, MAKGED, for knowledge graph error detection. Agents engage in several rounds of discussion and vote to determine triple correctness, thus enhancing the clarity and accuracy of decisions through collaborative evaluation.
    \item We design four trainable agents, including the ``\textit{Forward-Agent}'' and ``\textit{Backwa-
    rd-Agent}'' for the head entity and tail entity, respectively. These agents are trained on bidirectional subgraph embeddings, which enhances the robustness and adaptability of error detection.
    \item We build datasets with simulated graph noise, and experiments demonstrate that MAKGED improves accuracy by 0.73\% on FB15K and 6.62\% on WN18RR compared to the state-of-the-art methods. 
\end{enumerate}


\begin{figure*}[t]
  \centering
  \includegraphics[width=1\textwidth]{fig1.pdf}
  \caption{
  Training method for bidirectional subgraph agents in MAKGED: We first construct bidirectional subgraphs for the head and tail entities of the triple, which are represented as embedding vectors using GCN. These embeddings are then concatenated with the query embeddings of the Llama2 model, resulting in four bidirectional subgraph agents to evaluate the correctness of the triples.
  }
    \label{framework}
\end{figure*}


